\section{约束耦合的目标改善搜索}
\label{sec:coupled}

\subsection{将目标探测与可行性恢复结合}
非线性整数约束将多个变量的改变耦合在一起。仅改变一个目标相关坐标，可能立即破坏等式或乘积关系，使一个有希望的改善方向被拒绝。我们将目标探测和关联变量修复组织为联合动作：先提出希望保留的目标变化，再寻找能够恢复硬约束的配套坐标变化。

搜索维护探索赋值 $x$，而非强制每一步停留在当前最好赋值 $a$。候选先按加权未满足子句代价比较，再按修复距离和目标值比较。在可行区域内，普通动作要求保持硬约束并改善目标；离开可行域后，搜索以恢复约束为主，同时保留已经取得的严格改善方向。对于尚未优于 $b$ 的状态，虚拟严格改善惩罚使搜索避免反复返回已经取得的目标水平。

与这一评分相配合，私有修复问题可以显式加入
\begin{equation}
 f(x)-b+1\le0.
 \label{eq:strict-repair}
\end{equation}
由于目标为整数，式~\eqref{eq:strict-repair} 与 $f(x)<b$ 等价。修复副本保留全部硬约束，将改善目标转化为一次可行性搜索。该局部搜索的失败不能证明 $Q_b$ 不可满足。

\subsection{由约束依赖扩展单坐标动作}
设一次探测希望修改坐标集 $U$。投影过程保护这些请求坐标，尝试通过受影响的约束恢复其他坐标。对于含一次可解变量的等式，仅当代数修复具有整数解且位于搜索表示范围内时执行更新；对于可识别的常量条件定义，则根据守卫真值同步更新输出。投影只执行有限轮，不能修复的部分交由后续搜索承担。

例如，在 $x+y=4$ 下最小化 $x^2+y^2$。从 $(0,4)$ 单独改变一个坐标会违反等式，而联合动作 $(0,4)\to(1,3)$ 保持可行，并将目标从 $16$ 降为 $10$。这一例子说明，动作的评价需要考虑其可恢复的联合变化。

候选生成还包括联合非零因子修改、目标前沿探测和差分约束联动。差分修复将目标探测作为固定要求，寻找其余坐标的线性见证；在包含析取的情形下，若当前文字选择导致负环，则尝试替代的文字见证。这些过程均只产生候选，完整约束核验决定候选能否更新 $(a,b)$。

\subsection{因子覆盖下的线性化修复}
代数投影难以直接处理一般非线性关系。我们进一步选择一组固定因子，使其余变量可以在严格改善条件下共同调整。设硬约束和目标中的单项式为 $m(x)=\prod_i x_i^{e_i}$。固定坐标集 $Z$ 满足
\begin{equation}
 \forall m,\qquad \sum_{i\notin Z}e_i\le1.
 \label{eq:cover}
\end{equation}
给定固定坐标上的整数试点 $v$，得到限制问题
\begin{equation}
 F[x_Z:=v]\land f[x_Z:=v]<b.
 \label{eq:restricted-lia}
\end{equation}
式~\eqref{eq:cover} 同时覆盖硬约束和目标的全部单项式，因此式~\eqref{eq:restricted-lia} 属于线性整数算术（LIA）。对于二分双线性因子图，可通过颜色类构造固定集；更一般的多项式则选择满足次数条件的覆盖。试点包括当前赋值附近的值、小域组合、零值和较远临界点。

\begin{proposition}[覆盖线性化]
\label{prop:cover}
若 $Z$ 对硬约束和目标中的每个单项式均满足式~\eqref{eq:cover}，则任意整数 $v$ 下，式~\eqref{eq:restricted-lia} 的算术原子与目标均为仿射表达式。其满足赋值补上 $x_Z=v$ 后满足原 $F$，且目标小于 $b$。
\end{proposition}
\begin{proof}
固定变量的整数幂仅贡献常数系数，每项剩余总次数至多为一，故代入后的多项式均为仿射式。代入保留原 CNF 的布尔连接和全部原子，补回固定坐标即可恢复原问题的满足赋值及其严格改善目标。
\end{proof}

这一变换得到原问题的精确限制，而非过近似。线性局部搜索器在限制问题上允许其余坐标联合修复；返回赋值补齐固定坐标后，仍需在完整非线性问题上重新检查。有限试点失败或后端未找到解，均不构成最优性证明。这里的线性局部搜索用于寻找单个改善赋值，后文的方向局部搜索则寻找能够永久保持可行的整数方向，两者承担不同角色。

\begin{algorithm}[htbp]
\caption{\textsc{CoupledImprove}：联合改善搜索的一段工作}
\label{alg:coupled}
\begin{algorithmic}[1]
\Require $F,f,a,b$，探索状态 $\sigma$，软优先级 $H$，本片预算 $\tau$
\State 恢复 $\sigma$，独立保留可行参照 $a$
\While{本片预算未耗尽}
 \State 若 $\sigma$ 可行则优先选择目标相关坐标，否则选择冲突子句相关坐标
 \State 根据 $H$ 调整候选坐标的扫描顺序
 \State 生成临界点和目标探测，并对关联定义执行有限轮投影
 \State 按适用条件加入联合因子、差分联动和前沿候选
 \State 按可行性、修复代价和目标排序，执行选定动作
 \If{得到完整可行且严格改善的候选 $c$} \State \Return $\mathsf{Candidate}(c)$ \EndIf
 \State 若达到修复时机，则续搜严格界副本或覆盖线性化问题
 \If{修复得到完整可行改善候选 $c$} \State \Return $\mathsf{Candidate}(c)$ \EndIf
 \State 若停滞，则尝试中性移动、目标扰动或重启修复
\EndWhile
\State 保留 $\sigma$ 及可续传的修复状态；\Return $\mathsf{NoConclusion}$
\end{algorithmic}
\end{algorithm}

算法~\ref{alg:coupled} 将上述动作统一到一段可恢复搜索中。中性移动保持目标与可行性，并降低固定的约束压力；扰动则允许暂时离开可行域，但不会覆盖独立保存的 $a$。一次改善被接纳后，还可沿未标记为定义变量的整数坐标联合变化做倍增外推，每个点重新投影和检查。有限次外推是候选生成手段，不是永久可行方向的证明。
